\documentclass[12pt, letterpaper]{article}
\usepackage{graphicx} % Required for inserting images
%\usepackage[T1]{fontenc}

\title{Programy użytkowe \\ Lista 9}
\date{}
\usepackage{polski}
\begin{document}
\maketitle
\begin{enumerate}
    \item Zmodyfikuj zadanie z poprzedniej listy tak, aby sprawdzać, czy katalog już istnieje przed próbą jego utworzenia. Skorzystaj z z konstrukcji \verb|if[ -d ...]|. Uwaga na spacje wewnątrz nawiasu kwadratowego.
    \item Przy pomocy pętli \texttt{while read LINE} wczytaj zawartość pliku \texttt{wejście.txt} i wypisz na ekran 
    (\texttt{echo}) w następujący sposób
    \begin{enumerate}
        \item jeśli wartość w tym wierszu jest mniejsza niż 10000, zastąp ją przez 0;
        \item w przeciwnym wypadku, wypisz wiersz bez zmiany.
    \end{enumerate}
    \label{it:zadanie}
    \item Uruchom skrypt z poprzedniego zadania tak, aby wypisał wynik nie na ekran, ale do pliku \texttt{wyjście.txt}. Skorzystaj z operatora przekierowania \verb|>| (nie zmieniaj zawartości kodu).
    \item Zmodyfikuj skrypt z zadania \ref{it:zadanie} tak, aby wypisywanie zmodyfikowanych linijek odbywało się do pliku \texttt{wyjscie.txt}. Zwróć uwagę na różnicę między operatorami \verb|>| oraz \verb|>>|.
    
\end{enumerate}
\end{document}